% Generated draft: requires mathematical and source review.
\documentclass[11pt]{article}
\usepackage[margin=25mm]{geometry}
\usepackage{fontspec}
\setmainfont{Latin Modern Roman}
\usepackage{amsmath,amssymb,mathtools,unicode-math}
\setmathfont{Latin Modern Math}
\usepackage{enumitem,needspace,xurl,hyperref}
\setlength{\parindent}{0pt}
\setlength{\parskip}{5pt}
\setlength{\emergencystretch}{3em}
\newcommand{\N}{\mathbb{N}}
\newcommand{\Z}{\mathbb{Z}}
\newcommand{\Q}{\mathbb{Q}}
\newcommand{\R}{\mathbb{R}}
\newcommand{\C}{\mathbb{C}}
\newcommand{\F}{\mathbb{F}}
\newcommand{\A}{\mathbb{A}}
\newcommand{\PP}{\mathbb{P}}
\newcommand{\E}{\mathbb{E}}
\newcommand{\eps}{\varepsilon}
\newcommand{\dd}{\,\mathrm{d}}
\begin{document}
\section*{UnsolvedMath ID: 30003661}
\textbf{Source problem code:} catalogue entry
\subsection{Definitions and mathematical statement}
The supplied record does not specify the underlying represented space or the information supplied about closed sets. The following is a conditional formulation using the usual compact-square interpretation; these representation choices require confirmation against the source.

Let \(X=[0,1]^2\), with its Euclidean metric and the usual Cauchy representation of points. Represent a closed subset \(A\subseteq X\) by negative information: an enumeration of rational open balls, intersected with \(X\), whose union is \(X\setminus A\). Define
\[
\mathcal P_2=\{A\subseteq X:A\text{ is closed, nonempty, and pathwise connected}\}.
\]
Here pathwise connected means that for every \(x,y\in A\) there exists a continuous map \(\gamma:[0,1]\to A\) with \(\gamma(0)=x\) and \(\gamma(1)=y\). Pathwise connectedness is a promise on the input set; no paths or witnesses are supplied. Define the multivalued pathwise connected choice operation by
\[
\operatorname{PCC}_2:\mathcal P_2\rightrightarrows X,
\qquad \operatorname{PCC}_2(A)=A.
\]

Let \(2^{<\N}\) denote the finite binary strings and \(2^{\N}\) the infinite binary sequences. A binary tree is a prefix-closed subset \(T\subseteq 2^{<\N}\). Represent such trees by their characteristic functions, and restrict to trees having a node of every finite length. The multivalued operation associated with Weak K\"onig's Lemma is
\[
\operatorname{WKL}(T)=[T]
=\{p\in 2^{\N}:\forall n\in\N,\ p\mathbin{\upharpoonright}n\in T\},
\]
with the usual representation of binary sequences.

For multivalued operations \(f,g\) on represented spaces, write \(f\leq_{\mathrm{sW}}g\) if there exist computable partial transformations \(K,H\) of names such that, for every realizer \(G\) of \(g\), the composition \(H\circ G\circ K\) realizes \(f\). A realizer maps each valid input name to a name of an allowed output. In particular, \(K\) must produce a valid instance of \(g\), and \(H\) receives only the output name supplied by \(G\). Write \(f\equiv_{\mathrm{sW}}g\) when both strong reductions hold.

Under the conventions above, the question is whether
\[
\operatorname{PCC}_2\equiv_{\mathrm{sW}}\operatorname{WKL}.
\]
The record compares this question with the supplied assertion that connected choice is strongly Weihrauch equivalent to \(\operatorname{WKL}\) in every dimension at least \(2\). Under analogous cube and negative-information conventions, connected choice selects a point from a nonempty closed connected subset of \([0,1]^n\). The supplied literature assessment further reports that the corresponding geometric reduction for pathwise connected choice is known in dimensions at least \(3\), while dimension \(2\) remains open. These are supplied contextual claims, not independently verified results in this draft.
\subsection{Short English statement}
Given negative information describing a nonempty closed pathwise connected subset of the unit square, select a point in that set. Is this operation strongly Weihrauch equivalent to finding an infinite branch through an infinite binary tree? The square and representation conventions require source confirmation; the supplied record labels the question open.
\subsection{Sources}
\begin{itemize}
\item Supplied local registry record: permanent UnsolvedMath ID \(30003661\), title \emph{Pathwise Connected Choice in Dimension Two}, source problem code \texttt{catalogue entry}, in the Oberwolfach Reports Open Problems collection. No individual proposer is identified. The record gives proposed year \(2017\).
\item Supplied report citation: \emph{Mathematical Logic: Proof Theory, Constructive Mathematics} (2018), Oberwolfach Reports, pp.\ 3121--3183, DOI \texttt{10.4171/owr/2017/53}. \url{https://doi.org/10.4171/owr/2017/53}. The supplied background describes these reports as documenting mathematical workshops and their open-problem discussions.
\item Supplied literature reference, attributed to Brattka--Le Roux--Miller--Pauly: \url{https://people.math.wisc.edu/~jsmiller8/Papers/BFT.pdf}. No title, publication details, or theorem locator is supplied for this reference.
\item Literature assessment dated exactly \texttt{2026-08-21}: the record labels the dimension-two pathwise connected choice question open and reports the corresponding geometric reduction as known in dimensions three and higher. This assessment is supplied dated curation derived from the OWR material, not independently verified here; its status and cited claims require independent checking.
\item The registry labels the statement \texttt{corrected\_verified} and says its mathematical content was preserved while notation, typography, and sentence structure were repaired for publication. This records the supplied editorial assessment and does not establish independent mathematical verification of this draft.
\end{itemize}
\end{document}
